% Twin of ising_exact.md (md + tex + pdf; hand-written, 2026-11-13).  Keep in sync with the md
% by hand.  Compile from docs/: xelatex ising_exact.tex (twice), then delete aux/log/out.
\documentclass[11pt]{article}
\usepackage[margin=1in]{geometry}
\usepackage{amsmath,amssymb}
\usepackage{fontspec}
\setmonofont{Menlo}[Scale=MatchLowercase]
\usepackage{booktabs}
\usepackage{array}
\usepackage[hidelinks]{hyperref}
\setlength{\parindent}{0pt}
\setlength{\parskip}{0.6em}
\setlength{\emergencystretch}{3em}
\newcommand\code[1]{\texttt{#1}}
\DeclareMathOperator{\asinh}{asinh}
\DeclareMathOperator{\Tr}{Tr}
\title{The exact spectrum of the 2D Ising transfer matrix\\ on a strip of finite width}
\author{}\date{started 2026-11-13}
\begin{document}\maketitle

This note explains the formula behind \code{exact\_solutions/ising\_exact.py}: what object it
solves, the eigenvalue formula, the three conventions in use for saying which sign patterns are
eigenvalues, and how the momentum (spin) of a level is obtained.  The code is the reference
against which numerical transfer-matrix spectra are checked.

Every number quoted below was computed with that module on 2026-11-13.  The formula is not
derived here (section~3 sketches where it comes from and cites the derivation); what this note
can vouch for is that the formula \emph{as coded} reproduces the transfer matrix built from the
Boltzmann weights, level by level and label by label (section~8).

\section{The object}

Ising spins $s = \pm 1$ on a square lattice, coupling $K = J / k_B T$ on every bond.  Take $n$
spins in a row with periodic boundary conditions.  The row-to-row transfer matrix, in its
symmetric form, is
\begin{equation}
V[s, s'] = \exp\Big(\tfrac{K}{2} \sum_{i=1}^{n} s_i s_{i+1}\Big)\,
           \exp\Big(K \sum_{i=1}^{n} s_i s'_i\Big)\,
           \exp\Big(\tfrac{K}{2} \sum_{i=1}^{n} s'_i s'_{i+1}\Big),
\end{equation}
so that the partition function of the $n \times M$ torus is $Z = \Tr V^M$.  A ring of $n$
vertex tensors, or a ring of blocks after an untruncated coarse-graining, has the same spectrum.

$V$ commutes with two operations, so every eigenvector carries two labels:
\begin{itemize}
\item the \textbf{charge} $0$ or $1$: even or odd under the flip of all spins;
\item the \textbf{momentum} $P = 2\pi S / n$: translation by one site multiplies the eigenvector
  by $e^{iP}$.  The integer $S$, defined modulo $n$, is called the \textbf{spin} below (it is the
  conformal spin of the level at the critical point; it is not the Ising spin).
\end{itemize}

\section{The formula}

Define the dual coupling $K^*$ by
\begin{equation}
\tanh K^* = e^{-2K} \qquad\Longleftrightarrow\qquad \sinh 2K \, \sinh 2K^* = 1 ,
\end{equation}
and, for an angle $q$, the number $\gamma_q$ by
\begin{equation}
\cosh \gamma_q = \cosh 2K \, \cosh 2K^* - \cos q .
\end{equation}
Then every eigenvalue of $V$ has the form
\begin{equation}
\lambda = (2 \sinh 2K)^{n/2} \, \exp\Big( \tfrac12 \sum_{q} \epsilon_q \gamma_q \Big),
\qquad \epsilon_q = \pm 1 ,
\end{equation}
with one sign for each of the $n$ angles of its sector:
\begin{center}\begin{tabular}{lll}
\toprule
sector & angles & $j = 0, \dots, n-1$ \\
\midrule
charge 0 & $q = (2j+1)\,\pi / n$ & odd multiples of $\pi/n$ \\
charge 1 & $q = 2j\,\pi / n$ & even multiples of $\pi/n$ \\
\bottomrule
\end{tabular}\end{center}
Each sector has $2^n$ sign patterns and $2^{n-1}$ eigenvalues: half of the patterns occur.
Which half is the subject of section~4.

\paragraph{The dispersion in a form without cancellation.}  Because $\sinh 2K \sinh 2K^* = 1$,
one has $\cosh 2K \cosh 2K^* = 1 + \cosh 2(K - K^*)$, and the definition of $\gamma_q$ becomes
\begin{equation}
\sinh^2 \frac{\gamma_q}{2} = \sinh^2 (K - K^*) + \sin^2 \frac{q}{2}
\qquad\Longrightarrow\qquad
\gamma_q = 2 \asinh \sqrt{ \sinh^2 (K - K^*) + \sin^2 \tfrac{q}{2} } .
\end{equation}
This is the same function; the code uses it because the first form subtracts two numbers close
to 1 when $q$ is small or $T$ is near $T_c$.  Three special cases:
\begin{itemize}
\item $q = 0$: $\gamma_0 = 2\,|K - K^*|$.  The definition fixes $\gamma_0$ only up to its sign,
  which is where the conventions of section~4 come from.
\item $q = \pi$: $\gamma_\pi = 2\,(K + K^*)$, positive, no ambiguity.
\item at the critical point $K = K^* = \tfrac12 \ln(1 + \sqrt2)$:
  $\gamma_q = 2 \asinh |\sin \tfrac{q}{2}|$, which vanishes linearly at $q = 0$.
\end{itemize}
$\gamma_q = \gamma_{2\pi - q}$, so apart from $q = 0$ and $q = \pi$ the angles come in pairs of
equal $\gamma$.

\section{Where it comes from (a sketch)}

Write $V = V_1^{1/2} V_2 V_1^{1/2}$ with the bonds inside a row in
$V_1 = \exp\big(K \sum_i \sigma^z_i \sigma^z_{i+1}\big)$ and the bonds between two rows in $V_2$.
One vertical bond is the $2 \times 2$ matrix $e^{K} + e^{-K} \sigma^x$, and
\begin{equation}
e^{K} + e^{-K} \sigma^x = (2 \sinh 2K)^{1/2} \exp( K^* \sigma^x ),
\qquad\text{so}\qquad
V_2 = (2 \sinh 2K)^{n/2} \exp\Big( K^* \sum_i \sigma^x_i \Big) .
\end{equation}
That is where the prefactor and the dual coupling enter.  A Jordan--Wigner transformation turns
both exponents into quadratic forms in fermion operators.  The bond that closes the row picks up
the spin-flip parity, so the fermions are \emph{antiperiodic} around the strip in the charge-0
sector and \emph{periodic} in the charge-1 sector: hence the two sets of angles.  A Bogoliubov
rotation then diagonalizes the product mode by mode, and each mode contributes $e^{+\gamma_q/2}$
or $e^{-\gamma_q/2}$.

The mode $q = 0$ has no partner $-q$ to be rotated with.  Its two states, empty and occupied,
differ in energy by $2\,|K - K^*|$, and \emph{which} of the two is the lower one changes at
$T_c$.  That single fact is the origin of everything in section~4.  (The mode $q = \pi$, when
present, is unpaired too, but its energy $2(K + K^*)$ never changes sign.)

Derivations: L.~Onsager, Phys.\ Rev.\ \textbf{65}, 117 (1944); B.~Kaufman, Phys.\ Rev.\
\textbf{76}, 1232 (1949); T.~D.~Schultz, D.~C.~Mattis and E.~H.~Lieb, Rev.\ Mod.\ Phys.\
\textbf{36}, 856 (1964).

\section{Which sign patterns: three conventions, one set of levels}

In the charge-0 sector all codes agree: a pattern is an eigenvalue when the number of minus signs
is \emph{even}.  In the charge-1 sector they differ, and the whole difference sits in the one
mode $q = 0$, whose $\gamma_0$ the definition leaves undetermined in sign.

\begin{center}\small\renewcommand{\arraystretch}{1.15}%
\begin{tabular}{@{}>{\raggedright\arraybackslash}p{2.9cm}>{\raggedright\arraybackslash}p{2.7cm}>{\raggedright\arraybackslash}p{2.3cm}>{\raggedright\arraybackslash}p{2.8cm}>{\raggedright\arraybackslash}p{4.0cm}@{}}
\toprule
convention & value given to $\gamma_0$ & charge 0 & charge 1 & used by \\
\midrule
A, ``Kaufman'' & $2(K - K^*)$: positive below $T_c$ & even number of minus signs &
  \textbf{even} number of minus signs & \code{gammas}, \code{Ei\_kmax} of this module \\[2pt]
B, ``fermion vacuum'' & $2(K^* - K)$: positive above $T_c$ & even &
  \textbf{odd} number of minus signs & the usual fermionic presentation \\[2pt]
C, ``excitations'' & $2\,|K - K^*| \ge 0$ & even number of fermions &
  \textbf{odd above $T_c$, even below} & \code{levels}, \code{sector\_tables} of this module \\
\bottomrule
\end{tabular}\end{center}

Convention A is named after Kaufman because the script \code{gammas} and \code{Ei\_kmax} come
from follows her notation $H, H^*$ for $K, K^*$; her paper was not consulted for the sign.

\paragraph{Why A and B are the same set.}  Only the product $\epsilon_0 \gamma_0$ enters
$\lambda$.  Take a pattern allowed by B and reverse its first sign alone.  The product
$\epsilon_0 \gamma_0$ is unchanged, because the sign and the value given to $\gamma_0$ have both
been reversed, so $\lambda$ is the same number; and the count of minus signs has changed by
exactly one, so odd has become even: the new pattern is one allowed by A.  Doing it twice gives
back the original pattern, so the two lists are paired one to one, with the same multiplicities.

\paragraph{A worked example.}  Width $n = 3$, charge 1, at $1.2\,T_c$: $K = 0.36724$,
$K^* = 0.52267$, $\gamma_0 = \pm 0.31086$, $\gamma_{2\pi/3} = \gamma_{4\pi/3} = 1.58781$.  The
patterns are written $(\epsilon_0\, \epsilon_{2\pi/3}\, \epsilon_{4\pi/3})$.

\begin{center}\begin{tabular}{cccc}
\toprule
B: $\gamma_0 = +0.31086$, odd & A: $\gamma_0 = -0.31086$, even & $\lambda$ from either & $\lambda$ of $V$, diagonalized \\
\midrule
$(-++)$ & $(+++)$ & 8.5141788240 & 8.5141788240 \\
$(+-+)$ & $(--+)$ & 2.3744721434 & 2.3744721434 \\
$(++-)$ & $(-+-)$ & 2.3744721434 & 2.3744721434 \\
$(---)$ & $(+--)$ & 0.3556191894 & 0.3556191894 \\
\bottomrule
\end{tabular}\end{center}

Each row is one level written two ways; the two patterns differ in their first sign only.  The
last column is $V$ built from the Boltzmann weights of section~1 and diagonalized in the charge-1
sector.  Below $T_c$ the same four pairings hold; the numerical value of $\gamma_0$ has changed
sign in both conventions.

\paragraph{What each convention says.}
\begin{itemize}
\item B reads a minus sign as an occupied fermion, and the charge-1 sector always holds an odd
  number of them.  The $q = 0$ fermion costs $2(K^* - K)$, which is negative below $T_c$:
  occupying it there \emph{raises} $\lambda$.
\item A keeps one rule, ``even'', for both sectors and makes the largest eigenvalue of each
  sector the all-plus pattern.  The price is a $\gamma_0$ that changes sign at $T_c$.
\item C counts every $\gamma_q \ge 0$ as the energy of one fermion, so that more fermions always
  means a smaller $\lambda$.  Then the \emph{rule} has to change at $T_c$ instead of the sign: see
  the next section.
\end{itemize}
At $T_c$ itself $\gamma_0 = 0$, the first sign does nothing, and the three coincide.

\section{Levels as fermions}

Put $E = -\ln \lambda$.  In convention C
\begin{equation}
E = E_{\rm vac}({\rm charge}) + \sum_{q\ {\rm occupied}} \gamma_q ,
\qquad
E_{\rm vac} = -\frac{n}{2} \ln (2 \sinh 2K) - \frac12 \sum_{q} \gamma_q ,
\end{equation}
the sum in $E_{\rm vac}$ running over the $n$ angles of that sector, and the allowed occupations
are
\begin{itemize}
\item charge 0: an \textbf{even} number of fermions;
\item charge 1: an \textbf{odd} number above $T_c$, an \textbf{even} number below.
\end{itemize}
The lowest level of all, $E_0$, is the charge-0 vacuum.  The code reports
$X = (E - E_0)\, n / 2\pi$, which at $T_c$ tends to the scaling dimension as $n \to \infty$.

The change of rule at $T_c$ is the phase transition written in the counting.  Above $T_c$ the
lowest charge-1 level holds one fermion at $q = 0$, and its distance from $E_0$ tends to
$\gamma_0 = 2(K^* - K)$, the inverse correlation length.  Below $T_c$ it holds none and lies
exponentially close to $E_0$: the two ordered states.

\begin{center}\begin{tabular}{rll}
\toprule
$n$ & $E_1 - E_0$ at $1.2\,T_c$ \ ($\gamma_0 = 0.3108561287$) & $E_1 - E_0$ at $0.8\,T_c$ \\
\midrule
8 & 0.3261155415 & $7.7 \times 10^{-3}$ \\
16 & 0.3116924099 & $1.9 \times 10^{-4}$ \\
32 & 0.3108600721 & $1.8 \times 10^{-7}$ \\
64 & 0.3108561288 & $2.6 \times 10^{-13}$ \\
\bottomrule
\end{tabular}\end{center}

\paragraph{How many fermions the low levels hold.}  Nothing in \code{levels} caps the number of
fermions; levels come out in order of $E$.  At $T_c$ the cheapest $N$ fermions occupy the $N$
smallest $|q|$, which puts the first $N$-fermion level at $X \approx N^2/4$ in charge 0 ($N$
even) and at $X \approx \tfrac18 + (N^2 - 1)/4$ in charge 1 ($N$ odd).  Measured at $n = 1024$:

\begin{center}\begin{tabular}{rrrr}
\toprule
fermions & charge & first appears at $X$ & as level number \\
\midrule
1 & 1 & 0.125000 & 2 \\
2 & 0 & 0.999999 & 3 \\
3 & 1 & 2.124994 & 12 \\
4 & 0 & 3.999978 & 32 \\
5 & 1 & 6.124944 & 111 \\
6 & 0 & 8.999880 & 326 \\
8 & 0 & 15.999611 & 3,646 \\
10 & 0 & 24.999039 & 42,152 \\
\bottomrule
\end{tabular}\end{center}

So a list cut by particle number is complete only for the first few levels: a cap at two fermions
already fails at level 12, a cap at three at level 32.

\section{Momentum (spin)}

\paragraph{The rule.}  The momentum of a level is the sum of the momenta of its occupied
fermions, modulo $2\pi$; the vacuum of each sector has momentum zero.

\paragraph{In integers.}  Write every angle as $q = l\,\pi / n$ with an integer $l$ between $0$
and $2n - 1$: $l$ is odd in the charge-0 sector and even in the charge-1 sector.  A fermion at
that angle carries momentum $l/2$ in units of $2\pi/n$, and
\begin{equation}
S = \frac12 \sum_{\rm occupied} l \pmod n .
\end{equation}
The code carries $\sum l$ (its \code{twice\_spin}) rather than $S$ to avoid fractions.
\begin{itemize}
\item \textbf{$S$ is an integer.}  In the charge-0 sector a single fermion has half-integer
  momentum (odd $l$), which is the antiperiodic boundary condition of section~3; but the sector
  holds an even number of them, so $\sum l$ is even.  In the charge-1 sector every $l$ is even.
\item \textbf{The $q = 0$ fermion carries no momentum} ($l = 0$).  The spin of a level is
  therefore the same in conventions A, B and C, which differ only in how that one mode is counted.
\item \textbf{Negative momenta.}  The angle $2\pi - q$ is the momentum $-q$; it appears as
  $l \to 2n - l$.
\end{itemize}

\paragraph{The range.}  $S$ is defined modulo $n$.  The code reports the representative in
$[-n/2,\ n/2)$, computed as $((S + \lfloor n/2 \rfloor) \bmod n) - \lfloor n/2 \rfloor$.  For
$n = 8$ the labels run from $-4$ to $+3$: momentum $\pi$ is written $-4$ and $+4$ never appears.
For $n = 7$ they run from $-3$ to $+3$.

\paragraph{The sign is a convention.}  Counting the angles $0 < q < \pi$ as positive is a choice.
The opposite choice sends every $S$ to $-S$, and the spectrum is the same under that exchange, so
no comparison of eigenvalues can tell the two apart.  The label is not tied to a direction of
translation on the lattice; fixing that would need eigenvectors, which the module does not
produce.

\paragraph{A worked example.}  Width $n = 16$ at $T_c$, charge 0, two fermions.  In units where a
fermion at $l$ contributes $\gamma\, n / 2\pi$ to $X$: $l = 1$ and $l = 31$ each give $0.498401$,
$l = 3$ and $l = 29$ each give $1.458395$.

\begin{center}\begin{tabular}{ccrrl}
\toprule
occupied $l$ & momenta & $\sum l$ & $S$ & $X$ \\
\midrule
1, 31 & $+\tfrac12,\ -\tfrac12$ & 32 & 0 & 0.996803 \\
1, 3 & $+\tfrac12,\ +\tfrac32$ & 4 & $+2$ & 1.956796 \\
31, 29 & $-\tfrac12,\ -\tfrac32$ & 60 & $-2$ & 1.956796 \\
1, 29 & $+\tfrac12,\ -\tfrac32$ & 30 & $-1$ & 1.956796 \\
31, 3 & $-\tfrac12,\ +\tfrac32$ & 34 & $+1$ & 1.956796 \\
\bottomrule
\end{tabular}\end{center}

The first row is the energy operator.  The other four are one exactly four-fold degenerate level,
spread over four different spins.

\paragraph{The low levels at $T_c$,} $n = 64$, with the operator of the Ising CFT each one becomes
(the two with $S = \pm 2$ and $X \to 2$ are the stress tensor and its conjugate; which sign is
which is the convention above):

\begin{center}\begin{tabular}{lcccl}
\toprule
$X$ & charge & $S$ & fermions & becomes \\
\midrule
0.000000 & 0 & 0 & 0 & identity \\
0.125013 & 1 & 0 & 1 & $\sigma$, dimension $1/8$ \\
0.999799 & 0 & 0 & 2 & $\varepsilon$, dimension $1$ \\
1.124210 & 1 & $\pm 1$ & 1 & derivatives of $\sigma$ \\
1.997196 & 0 & $\pm 1$ & 2 & derivatives of $\varepsilon$ \\
1.997196 & 0 & $\pm 2$ & 2 & $T$, $\bar T$ \\
2.118618 & 1 & $\pm 2$ & 1 & second-level descendants of $\sigma$ \\
2.123408 & 1 & 0 & 3 & second-level descendant of $\sigma$ \\
\bottomrule
\end{tabular}\end{center}

The ground level also carries the central charge: $E_0 = n f - \pi c / 6n + \dots$ with Onsager's
bulk free energy $f$ gives $c = 1/2$ to $2 \times 10^{-4}$ at $n = 128$.

\section{Comparing with a numerical transfer matrix}

\begin{itemize}
\item \textbf{Ask by number of levels, per sector.}  A solver returns the lowest $N$ eigenvalues
  of a sector, in order, and knows nothing of fermion number.  \code{sector\_tables(n, T, N)} and
  \code{levels(n, T, N, charge=, spin=)} return the same thing, complete.
\item \textbf{$n$ is the full width in lattice sites}: for a ring of $r$ blocks of $L$ sites,
  $n = rL$.
\item \textbf{Momentum is finer on the exact side.}  A ring of $r$ blocks, translated block by
  block, resolves the spin only modulo $r$; group the exact levels by $S \bmod r$ before
  comparing.  Treat $+S$ and $-S$ as a pair.
\item \textbf{Do not compare one global list across a degenerate level.}  A Krylov solver started
  from one vector can return two copies of a four-fold level and then the next level.  Sector by
  sector the four copies are in four different lists and the problem does not arise.
\end{itemize}

\section{The code, and how it was checked}

\begin{center}\begin{tabular}{ll}
\toprule
in this note & in \code{exact\_solutions/ising\_exact.py} \\
\midrule
$K^*$, $K - K^*$, $\gamma_q$ & \code{dual\_coupling}, \code{detuning}, \code{dispersion} \\
the $2n$ values $\gamma_l$, convention A & \code{gammas(T, n)} \\
the lowest levels, convention A, no labels & \code{Ei\_kmax(n, T, num\_collect, kmax=None)} \\
the lowest levels with charge, $S$, fermions & \code{levels(n, T, num\_levels, charge=, spin=)} \\
the lowest levels of all $2n$ sectors & \code{sector\_tables(n, T, num\_levels)} \\
$E_0$ & \code{ground\_energy(n, T)} \\
\bottomrule
\end{tabular}\end{center}

\code{levels} finds the lowest $N$ of the whole spectrum best-first with a heap;
\code{sector\_tables} finds the lowest $N$ of every sector by a dynamic programme over the modes,
whose cost does not depend on the spin.  The \code{README.md} has the timings and the
provenance.

Checks (\code{tests/test\_ising\_exact.py}, 13 tests):
\begin{itemize}
\item \textbf{Against the definition.}  $V$ of section~1 built from the Boltzmann weights and
  projected onto every (charge, momentum) sector, with the spin flip and the translation written
  as permutations of the $2^n$ configurations: all $2^n$ levels, their values and both labels,
  for $n = 1, \dots, 8$ at $0.7$, $1$ and $1.4\,T_c$, to $10^{-10}$, by both algorithms.
  Shifting the charge-1 spins by one, or exchanging the two charges, makes the comparison fail,
  so it does test the labels; it cannot test the sign of $S$.  It is also the evidence that both
  vacua have momentum zero.
\item \textbf{Against the recorded values} that the earlier script returned, kept in
  \code{tests/\allowbreak{}data/\allowbreak{}original\_reference.json}: 900 argument
  combinations of \code{Ei\_kmax}, to $10^{-11}$.
\item \textbf{The two algorithms against each other} at $n = 24$, and the CFT content and
  $c = 1/2$ above.
\end{itemize}

\end{document}
